\subsection{Explanation}
Each subsection is an aspect of the project (either the documentation, Python, or Java application). Each subsubsection is then an epic within that area. Each paragraph represents one story for that epic, referencing a story in the project backlog. Each story is assigned a value associated with how much work is required to complete the story for this sprint, the risks associated with the story, the technical changes required to complete it, what it depends on, how complex it is to achieve, and what value it adds to the project, assuming all work to be done in the last sprint was done. Each story is also assigned a priority of either must, should, could, or wont. Lastly, an acceptance criteria is provided for each story, stating what tasks must be completed to achieve the story.

\subsection{Documentation}

\subsubsection{Report}

\paragraph{Stakeholders Can View Executive Summary}
\textbf{Value:} 2 \quad \textbf{Priority:} Must\\
\textbf{Acceptance criteria}
\begin{itemize}
  \item The summary appears at the start of the report.
  \item The summary clearly states what is delivered in the prototype.
  \item The summary clearly states why features exist in the prototype.
  \item The summary clearly states the most important risks of the prototype.
  \item The summary clearly states the mitigations for those risks.
\end{itemize}

\paragraph{Stakeholders Can View Problem Framing}
\textbf{Value:} 3\quad \textbf{Priority:} Must\\
\textbf{Acceptance criteria}
\begin{itemize}
  \item The problem framing clearly explains the problem.
  \item The problem framing clearly states why the problem is important.
  \item The problem framing clearly states what the goal of the project is.
  \item The problem framing states who is affected by the problem.
\end{itemize}

\paragraph{Stakeholders Can View Definition of Done}
\textbf{Value:} 1\quad \textbf{Priority:} Must\\
\textbf{Acceptance criteria}
\begin{itemize}
  \item The definition of done clearly explains when a documentation story is complete.
  \item The definition of done clearly explains when a program story is complete.
\end{itemize}

\paragraph{Stakeholders Can View Project Backlog}
\textbf{Value:} 3\quad \textbf{Priority:} Must\\
\textbf{Acceptance criteria}
\begin{itemize}
  \item The project backlog clearly states which epics are required to complete the project.
  \item The project backlog clearly states which user stories are associated with each epic.
\end{itemize}

\paragraph{Stakeholders Can View Sprint One Prioritised Requirements}
\textbf{Value:} 3\quad \textbf{Priority:} Must\\
\textbf{Acceptance criteria}
\begin{itemize}
  \item Each epic in the project backlog is stated in the requirements.
  \item Each story within each epic is stated in the requirements.
  \item Each story has a value associated with it, based on how much work it is to complete the story, the significance of the risks associated with it, the technical changes required to complete it, how much depends on it, how complex it is to complete, and the value its completion adds to the project.
  \item Each story has an acceptance criteria associated with it, which, when all criteria are met, is considered complete for sprint one. Criteria may change between sprints as sprint requirements change.
  \item Each story has a priority associated with it for that sprint. The priority must be must, should, could, or wont.
  \item Each story is able to link to a prototype feature.
\end{itemize}

\paragraph{Stakeholders Can View Architecture Schema}
\textbf{Value:} 3\quad \textbf{Priority:} Should\\
\textbf{Acceptance criteria}
\begin{itemize}
  \item Each class category of the Java application has its purpose and use case explained.
  \item Each class category of the Python application has its purpose and use case explained.
  \item Classes that share an identical structure with minor variations need only have their category explained.
  \item All structural design choices are justified and explained.
  \item The limitations of all specified design choices are stated, and why they are acceptable is explained.
\end{itemize}

\paragraph{Stakeholders Can View Data Flow Schema}
\textbf{Value:} 2\quad \textbf{Priority:} Should\\
\textbf{Acceptance criteria}
\begin{itemize}
  \item Each component that accesses data is explained, describing what data it accesses, what it does with it, and why it needs to access this data.
  \item Each data store is explained, stating what data it stores and in what format.
\end{itemize}

\paragraph{Stakeholders Can View Telemetry Schema}
\textbf{Value:} 3\quad \textbf{Priority:} Must\\
\textbf{Acceptance criteria}
\begin{itemize}
  \item Each telemetry event is described by the schema.
  \item Each telemetry event has all its fields described by the schema.
  \item Each field has its domain described by the schema. For enumerated fields, all valid values are described.
\end{itemize}

\paragraph{Stakeholders Can View Game Design Documentation}
\textbf{Value:} 3\quad \textbf{Priority:} Could\\
\textbf{Acceptance criteria}
\begin{itemize}
  \item The documentation outlines the gameplay loop.
  \item The documentation explains all parts of the gameplay loop.
  \item The documentation defines all attack abilities (both magic and mundane) along with their cost to purchase, magic cost, damage type, and base damage.
  \item The documentation defines all passive abilities, what they do, and their cost.
  \item The documentation defines all damage types and their use cases.
  \item The documentation outlines a structure for the game's progression.
  \item The documentation defines all enemies within the game.
\end{itemize}

\paragraph{Stakeholders Can View Evaluation}
\textbf{Value:} 3\quad \textbf{Priority:} Must\\
\textbf{Acceptance criteria}
\begin{itemize}
  \item The evaluation defines five measures of success.
  \item For each measure, evidence is provided to show system effectiveness.
  \item The limitations and biases of the evaluation are explored and explained.
\end{itemize}

\paragraph{Stakeholders Can View Sprint Two Prioritised Requirements}
\textbf{Value:} 3\quad \textbf{Priority:} Should\\
\textbf{Acceptance criteria}
\begin{itemize}
  \item Each epic in the project backlog is stated in the requirements.
  \item Each story within each epic is stated in the requirements.
  \item Each story has a value associated with it, based on how much work it is to complete the story, the significance of the risks associated with it, the technical changes required to complete it, how much depends on it, how complex it is to complete, and the value its completion adds to the project.
  \item Each story has an acceptance criteria associated with it, which, when all criteria are met, is considered complete for sprint two. Criteria may change between sprints as sprint requirements change.
  \item Each story has a priority associated with it for that sprint. The priority must be must, should, could, or wont.
  \item Each story is able to link to a final product feature and test.
\end{itemize}

\paragraph{Stakeholders Can View Management}
\textbf{Value:} 2\quad \textbf{Priority:} Should\\
\textbf{Acceptance criteria}
\begin{itemize}
  \item The management section outlines what roles each member had.
  \item It outlines what tasks members owned and what tasks they contributed to.
  \item It outlines the areas in which the group worked well together.
  \item It outlines the challenges the group faced working together.
  \item It outlines how those challenges were solved.
\end{itemize}

\subsubsection{Handover Pack}

\paragraph{Client Can View Application Deployment and Operation Instructions}
\textbf{Value:} 1\quad \textbf{Priority:} Must\\
\textbf{Acceptance criteria}
\begin{itemize}
  \item The guide contains clear instructions for how to set up and operate both applications.
  \item There may be instructions for different systems, but there must be a way to run the applications on any desktop system that runs Windows, macOS, or Linux.
\end{itemize}

\paragraph{Maintainer Can View the Automated Test Running Instructions}
\textbf{Value:} 1\quad \textbf{Priority:} Must\\
\textbf{Acceptance criteria}
\begin{itemize}
  \item There are clear instructions on how to operate the test suite on the source code.
  \item The version of the testing framework is provided.
  \item The running instructions allow the tests to be run on any desktop system that runs Windows, macOS, or Linux.
\end{itemize}

\paragraph{Maintainers Can View Maintenance Instructions}
\textbf{Value:} 0\quad \textbf{Priority:} Wont\\
\textbf{Acceptance criteria} N/A

\paragraph{Users Can View Troubleshooting Instructions}
\textbf{Value:} 0\quad \textbf{Priority:} Wont\\
\textbf{Acceptance criteria} N/A

\paragraph{Maintainers Can View Extension Instructions}
\textbf{Value:} 0\quad \textbf{Priority:} Wont\\
\textbf{Acceptance criteria} N/A

\paragraph{Maintainers Can View Data Management Instructions}
\textbf{Value:} 0\quad \textbf{Priority:} Wont\\
\textbf{Acceptance criteria} N/A

\subsubsection{Process Information}

\paragraph{Stakeholders Can View Meeting Minutes}
\textbf{Value:} 3\quad \textbf{Priority:} Must\\
\textbf{Acceptance criteria}
\begin{itemize}
  \item The meeting minutes state who attended each meeting.
  \item They state what work was completed by each group member who attended.
  \item They state what topics were discussed in each meeting.
  \item They state what work is to be done in the next weekly scrum and by whom.
\end{itemize}

\paragraph{Stakeholder Can View Risk Register}
\textbf{Value:} 5\quad \textbf{Priority:} Must\\
\textbf{Acceptance criteria}
\begin{itemize}
  \item Each risk is assigned a unique identifier.
  \item Each risk is given a clear description.
  \item The consequences of each risk occurring are clearly stated.
  \item The ordinal impact, from 1 to 5, is stated for each risk.
  \item The ordinal likelihood, from 1 to 5, is stated for each risk.
  \item The ordinal value, from 1 to 25, equalling the product of the impact and likelihood, is stated.
  \item The treatment for each risk is stated, including treatment method (accept, transfer, mitigate, avoid) and details of the treatment.
  \item An owner for each risk is assigned and recorded.
\end{itemize}

\paragraph{Stakeholder Can View Scrum Board}
\textbf{Value:} 3\quad \textbf{Priority:} Must\\
\textbf{Acceptance criteria}
\begin{itemize}
  \item There is a card for each story in the project backlog defined in this report.
  \item There is a card for each group of acceptance criteria (task) that varies between sprints.
  \item Epics are created on the board for each epic defined in the project backlog in this report.
  \item There is a backlog section for tasks that have yet to be started.
  \item There is an in-progress section for tasks being worked on. Each task in this section is assigned to at least one group member.
  \item There is a done section for tasks that are complete. Each task in this section is assigned to at least one group member.
\end{itemize}

\subsubsection{Ethics and Licensing}

\paragraph{Stakeholders Can View the Privacy and Data Protection Analysis}
\textbf{Value:} 1\quad \textbf{Priority:} Must\\
\textbf{Acceptance criteria}
\begin{itemize}
  \item It discusses what data is stored about users.
  \item It discusses how the data storage complies with law.
  \item It discusses how collected data is pseudonymised.
  \item It discusses how access to this data is limited to only those it is relevant to.
\end{itemize}

\paragraph{Stakeholders Can View the Consent and Disclosure Analysis}
\textbf{Value:} 1\quad \textbf{Priority:} Must\\
\textbf{Acceptance criteria}
\begin{itemize}
  \item It discusses what the data is used for.
  \item It discusses how consent is handled.
  \item It discusses how the consent and disclosure policy complies with law.
  \item It discusses how this policy affects the user experience.
\end{itemize}

\paragraph{Stakeholders Can View the Accessibility Analysis}
\textbf{Value:} 1\quad \textbf{Priority:} Should\\
\textbf{Acceptance criteria}
\begin{itemize}
  \item It discusses the importance of accessibility.
  \item It discusses what systems are in place to support the accessibility of the software.
  \item It discusses what areas the system can improve upon for sprint two.
\end{itemize}

\paragraph{Stakeholders Can View the Intellectual Property and Licensing Implications}
\textbf{Value:} 1\quad \textbf{Priority:} Must\\
\textbf{Acceptance criteria}
\begin{itemize}
  \item It discusses the intellectual property rights of the developers.
  \item It discusses the intellectual property rights of the product owner.
  \item It discusses the licensing considerations for the software and data the system depends on.
  \item It discusses the intellectual property rights of sources the system draws inspiration from or is similar to.
\end{itemize}

\paragraph{Stakeholders Can View the Project License}
\textbf{Value:} 1\quad \textbf{Priority:} Must\\
\textbf{Acceptance criteria}
\begin{itemize}
  \item A license for the project is determined.
  \item The license complies with the licensing requirements of all dependencies.
  \item The reasons for the choice of license are explained.
\end{itemize}

\paragraph{Maintainers Can View the Software and Data Inventory}
\textbf{Value:} 2\quad \textbf{Priority:} Should\\
\textbf{Acceptance criteria}
\begin{itemize}
  \item An inventory of all software and data components the system directly depends on is produced.
  \item It contains the license of each dependency.
  \item It contains the cost model of each dependency.
  \item It contains the provenance of each dependency.
  \item It contains the version of each dependency the system uses.
\end{itemize}

\subsubsection{Technical Information}

\paragraph{Stakeholder Can View the Automated Test Evidence}
\textbf{Value:} 2\quad \textbf{Priority:} Must\\
\textbf{Acceptance criteria}
\begin{itemize}
  \item There are five automated tests.
  \item Some of them are unit tests.
  \item Others are integration tests.
\end{itemize}

\paragraph{Stakeholder Can View the End-to-End Test Evidence}
\textbf{Value:} 1\quad \textbf{Priority:} Must\\
\textbf{Acceptance criteria}
\begin{itemize}
  \item There is one end-to-end test.
  \item Evidence that the test is successful is given.
  \item The happy path and failure cases for the test are described.
  \item The expected results of the test are described.
\end{itemize}

\paragraph{Stakeholders Can View Security Checklist}
\textbf{Value:} 0\quad \textbf{Priority:} Wont\\
\textbf{Acceptance criteria} N/A

\subsubsection{Presentation}

\paragraph{Stakeholder Can View an Explanation of the Problem}
\textbf{Value:} 2\quad \textbf{Priority:} Must\\
\textbf{Acceptance criteria}
\begin{itemize}
  \item It outlines what the problem is.
  \item It outlines why the problem is relevant.
  \item It outlines the scope of the problem.
  \item It outlines the scope of the solution.
  \item It outlines what assumptions were made about the problem.
  \item It outlines who the intended users of the solution are.
\end{itemize}

\paragraph{Stakeholder Can View the Design Approach}
\textbf{Value:} 2\quad \textbf{Priority:} Must\\
\textbf{Acceptance criteria}
\begin{itemize}
  \item It outlines the architecture of the solution.
  \item It outlines measures used to evaluate the system.
\end{itemize}

\paragraph{Stakeholder Can View Details of the Implementation}
\textbf{Value:} 2\quad \textbf{Priority:} Must\\
\textbf{Acceptance criteria}
\begin{itemize}
  \item It outlines the key components of the solution.
  \item It outlines the key features of the implementation.
  \item It outlines why key design decisions were made.
\end{itemize}

\paragraph{Stakeholder Can View the Evaluation}
\textbf{Value:} 2\quad \textbf{Priority:} Must\\
\textbf{Acceptance criteria}
\begin{itemize}
  \item It outlines what issues were faced during implementation.
  \item It outlines how well the system performs against five measures, with clear evidence.
\end{itemize}

\paragraph{Stakeholder Can View a Live Demonstration}
\textbf{Value:} 3\quad \textbf{Priority:} Must\\
\textbf{Acceptance criteria}
\begin{itemize}
  \item It takes around four minutes to complete.
  \item It demonstrates a run of the game.
  \item It demonstrates the telemetry views.
  \item It demonstrates the telemetry suggestions.
  \item It demonstrates the simulation.
  \item It demonstrates the adjustment of a design parameter.
  \item It is rehearsed before the live demo to ensure it works correctly.
  \item It has a backup demonstration if it goes wrong.
\end{itemize}

\paragraph{Stakeholder Can View the Limitations and Next Steps}
\textbf{Value:} 2\quad \textbf{Priority:} Must\\
\textbf{Acceptance criteria}
\begin{itemize}
  \item It outlines the areas in which the prototype is weakest.
  \item It briefly outlines what is to be done in the next sprint.
\end{itemize}

\paragraph{Stakeholder Can View the Ethics and Telemetry Disclosure}
\textbf{Value:} 0\quad \textbf{Priority:} Wont\\
\textbf{Acceptance criteria} N/A

\paragraph{Stakeholders Can View Management Information}
\textbf{Value:} 1\quad \textbf{Priority:} Must\\
\textbf{Acceptance criteria}
\begin{itemize}
  \item It outlines the role of each member.
  \item It outlines what each member contributed to the project.
\end{itemize}

\paragraph{Stakeholders Can View a Brief of the Scrum Board}
\textbf{Value:} 0\quad \textbf{Priority:} Wont\\
\textbf{Acceptance criteria} N/A

\subsection{Python Application}

\subsubsection{Authentication and Access Control}

\paragraph{Users Can Log In}
\textbf{Value:} 5\quad \textbf{Priority:} Must\\
\textbf{Acceptance criteria}
\begin{itemize}
  \item Users must be required to log in before accessing the Python application.
  \item Users with an account must be able to log into that account by entering the correct details.
  \item Users without an account must be able to create one.
\end{itemize}

\paragraph{Users Can Reset Their Login}
\textbf{Value:} 1\quad \textbf{Priority:} Could\\
\textbf{Acceptance criteria}
\begin{itemize}
  \item Users must be able to reset their login credentials.
  \item This reset system must require them to have some other security factor such as email.
\end{itemize}

\paragraph{Users Can Perform Actions Permitted by Their Role}
\textbf{Value:} 2\quad \textbf{Priority:} Must\\
\textbf{Acceptance criteria}
\begin{itemize}
  \item A user with the player role must not be able to access the application.
  \item A user with the designer or developer role must have full access to the application.
\end{itemize}

\paragraph{Users Can View Their Profile}
\textbf{Value:} 0\quad \textbf{Priority:} Wont\\
\textbf{Acceptance criteria} N/A

\subsubsection{Telemetry Views}

\paragraph{Designers Can View Funnel}
\textbf{Value:} 3\quad \textbf{Priority:} Must\\
\textbf{Acceptance criteria}
\begin{itemize}
  \item The system calculates a funnel view showing user character drop-off from the telemetry data as the game progresses.
  \item This is then displayed to the user.
  \item It is able to be switched from viewing actual user data to simulation generated telemetry data.
\end{itemize}

\paragraph{Designers Can View Difficulty Spikes}
\textbf{Value:} 3\quad \textbf{Priority:} Must\\
\textbf{Acceptance criteria}
\begin{itemize}
  \item The system calculates a difficulty spikes view showing failure rate for different stages from the telemetry data.
  \item This is then displayed to the user.
  \item It is able to be switched from viewing actual user data to simulation generated telemetry data.
\end{itemize}

\paragraph{Designers Can View Progress Curves}
\textbf{Value:} 3 \quad \textbf{Priority:} Could\\
\textbf{Acceptance criteria}
\begin{itemize}
  \item The system calculates how many coins users' characters accumulate on each stage from the telemetry data.
  \item This is then displayed to the user.
  \item It is able to be switched from viewing actual user data to simulation generated telemetry data.
\end{itemize}

\paragraph{Designers Can View Fairness Indicators}
\textbf{Value:} 0 \quad \textbf{Priority:} Wont\\
\textbf{Acceptance criteria} N/A

\paragraph{Designers Can View Comparison Mode}
\textbf{Value:} 3 \quad \textbf{Priority:} Could\\
\textbf{Acceptance criteria}
\begin{itemize}
  \item The system calculates how many coins users' characters accumulate on each stage, along with how much health they finish each stage with, for each difficulty.
  \item This is then displayed to the user.
  \item It is able to be switched from viewing actual user data to simulation generated telemetry data.
\end{itemize}

\paragraph{Designers Can View Rule Based Design Suggestions}
\textbf{Value:} 2 \quad \textbf{Priority:} Should\\
\textbf{Acceptance criteria}
\begin{itemize}
  \item The system reads the telemetry data and evaluates it against a suggestion criteria.
  \item If this criteria is met, the system suggests a concrete change to the design parameter to improve the balance of the game.
  \item These suggestions are displayed to the user.
\end{itemize}

\paragraph{Designers Can View Decision Log}
\textbf{Value:} 0 \quad \textbf{Priority:} Wont\\
\textbf{Acceptance criteria} N/A

\subsubsection{Telemetry Events}

\paragraph{Application Handles Start Session}
\textbf{Value:} 1 \quad \textbf{Priority:} Must\\
\textbf{Acceptance criteria}
\begin{itemize}
  \item The application is able to read start session events that follow the telemetry schema.
  \item The application is able to validate those events.
  \item The application provides clear errors to the user when failing to validate them.
\end{itemize}

\paragraph{Application Handles Normal Encounter Fail}
\textbf{Value:} 1 \quad \textbf{Priority:} Must\\
\textbf{Acceptance criteria}
\begin{itemize}
  \item The application is able to read normal encounter fail events that follow the telemetry schema.
  \item The application is able to validate those events.
  \item The application provides clear errors to the user when failing to validate them.
\end{itemize}

\paragraph{Application Handles Normal Encounter Complete}
\textbf{Value:} 1 \quad \textbf{Priority:} Must\\
\textbf{Acceptance criteria}
\begin{itemize}
  \item The application is able to read normal encounter complete events that follow the telemetry schema.
  \item The application is able to validate those events.
  \item The application provides clear errors to the user when failing to validate them.
\end{itemize}

\paragraph{Application Handles Boss Encounter Start}
\textbf{Value:} 1 \quad \textbf{Priority:} Could\\
\textbf{Acceptance criteria}
\begin{itemize}
  \item The application is able to read boss encounter start events that follow the telemetry schema.
  \item The application is able to validate those events.
  \item The application provides clear errors to the user when failing to validate them.
\end{itemize}

\paragraph{Application Handles Boss Encounter Fail}
\textbf{Value:} 1 \quad \textbf{Priority:} Could\\
\textbf{Acceptance criteria}
\begin{itemize}
  \item The application is able to read boss encounter fail events that follow the telemetry schema.
  \item The application is able to validate those events.
  \item The application provides clear errors to the user when failing to validate them.
\end{itemize}

\paragraph{Application Handles Boss Encounter Complete}
\textbf{Value:} 1 \quad \textbf{Priority:} Could\\
\textbf{Acceptance criteria}
\begin{itemize}
  \item The application is able to read boss encounter complete events that follow the telemetry schema.
  \item The application is able to validate those events.
  \item The application provides clear errors to the user when failing to validate them.
\end{itemize}

\paragraph{Application Handles Gain Coin}
\textbf{Value:} 1 \quad \textbf{Priority:} Could\\
\textbf{Acceptance criteria}
\begin{itemize}
  \item The application is able to read gain coin events that follow the telemetry schema.
  \item The application is able to validate those events.
  \item The application provides clear errors to the user when failing to validate them.
\end{itemize}

\paragraph{Application Handles Buy Upgrade}
\textbf{Value:} 1 \quad \textbf{Priority:} Could\\
\textbf{Acceptance criteria}
\begin{itemize}
  \item The application is able to read buy upgrade events that follow the telemetry schema.
  \item The application is able to validate those events.
  \item The application provides clear errors to the user when failing to validate them.
\end{itemize}

\paragraph{Application Handles End Session}
\textbf{Value:} 1 \quad \textbf{Priority:} Must\\
\textbf{Acceptance criteria}
\begin{itemize}
  \item The application is able to read end session events that follow the telemetry schema.
  \item The application is able to validate those events.
  \item The application provides clear errors to the user when failing to validate them.
\end{itemize}

\paragraph{Application Handles Settings Change}
\textbf{Value:} 1 \quad \textbf{Priority:} Must\\
\textbf{Acceptance criteria}
\begin{itemize}
  \item The application is able to read settings change events that follow the telemetry schema.
  \item The application is able to validate those events.
  \item The application provides clear errors to the user when failing to validate them.
\end{itemize}

\paragraph{Application Handles Kill Enemy}
\textbf{Value:} 1 \quad \textbf{Priority:} Could\\
\textbf{Acceptance criteria}
\begin{itemize}
  \item The application is able to read kill enemy events that follow the telemetry schema.
  \item The application is able to validate those events.
  \item The application provides clear errors to the user when failing to validate them.
\end{itemize}

\subsubsection{Seeded Dataset}

\paragraph{Developers Can Access a Telemetry Dataset}
\textbf{Value:} 0 \quad \textbf{Priority:} Wont\\
\textbf{Acceptance criteria} N/A

\paragraph{Developers Can Access a Decision Dataset}
\textbf{Value:} 0 \quad \textbf{Priority:} Wont\\
\textbf{Acceptance criteria} N/A

\subsection{Java Application}

\subsubsection{Authentication and Access Control}

\paragraph{User Can Log In}
\textbf{Value:} 5\quad \textbf{Priority:} Must\\
\textbf{Acceptance criteria}
\begin{itemize}
  \item Users must be required to log in before accessing the Java application.
  \item Users with an account must be able to log into that account by entering the correct details.
  \item Users without an account must be able to create one.
\end{itemize}

\paragraph{Users Can Reset Their Login}
\textbf{Value:} 1\quad \textbf{Priority:} Could\\
\textbf{Acceptance criteria}
\begin{itemize}
  \item Users must be able to reset their login credentials.
  \item This reset system must require them to have some other security factor such as email.
\end{itemize}

\paragraph{Users Can Perform Actions Permitted by Their Role}
\textbf{Value:} 2\quad \textbf{Priority:} Must\\
\textbf{Acceptance criteria}
\begin{itemize}
  \item A user with the player role is able to toggle telemetry, view the telemetry disclosure, and play the game.
  \item A user with the designer or developer role is able to view the values of design parameters and execute simulated runs.
\end{itemize}

\subsubsection{Settings}

\paragraph{Users Can View the Telemetry Disclosure}
\textbf{Value:} 1\quad \textbf{Priority:} Must\\
\textbf{Acceptance criteria}
\begin{itemize}
  \item A user in the main menu is able to view the telemetry disclosure.
  \item A user in settings is able to view the telemetry disclosure.
  \item The telemetry disclosure contains what data is collected.
  \item The telemetry disclosure explains why that data is collected.
\end{itemize}

\paragraph{Users Can Toggle Telemetry}
\textbf{Value:} 1\quad \textbf{Priority:} Must\\
\textbf{Acceptance criteria}
\begin{itemize}
  \item A user in settings is able to see if they have telemetry enabled or not.
  \item A user in settings is able to toggle whether they have telemetry enabled or not.
  \item When disabled, telemetry events must not be recorded.
  \item It must, by default, be on.
  \item The setting must persist beyond the lifetime of the game.
\end{itemize}

\paragraph{Designers Can View and Change Starting Lives}
\textbf{Value:} 2\quad \textbf{Priority:} Must\\
\textbf{Acceptance criteria}
\begin{itemize}
  \item A designer in settings is able to view the starting lives for each difficulty.
  \item A designer in settings is able to set the value of starting lives for each difficulty.
  \item The value for starting lives persists beyond the lifetime of the game.
\end{itemize}

\paragraph{Designers Can View and Change Enemy Health Multiplier}
\textbf{Value:} 0\quad \textbf{Priority:} Wont\\
\textbf{Acceptance criteria} N/A

\paragraph{Designers Can View and Change Player Maximum Health}
\textbf{Value:} 0\quad \textbf{Priority:} Wont\\
\textbf{Acceptance criteria} N/A

\paragraph{Designers Can View and Change Upgrade Price Multiplier}
\textbf{Value:} 0\quad \textbf{Priority:} Wont\\
\textbf{Acceptance criteria} N/A

\paragraph{Designers Can View and Change Enemy Damage Multiplier}
\textbf{Value:} 0\quad \textbf{Priority:} Wont\\
\textbf{Acceptance criteria} N/A

\paragraph{Designers Can View and Change Player Maximum Magic}
\textbf{Value:} 0\quad \textbf{Priority:} Wont\\
\textbf{Acceptance criteria} N/A

\paragraph{Designers Can View and Change Magic Regeneration Rate}
\textbf{Value:} 0\quad \textbf{Priority:} Wont\\
\textbf{Acceptance criteria} N/A

\paragraph{Designers Can View and Change Number of Items in the Shop}
\textbf{Value:} 0\quad \textbf{Priority:} Wont\\
\textbf{Acceptance criteria} N/A

\paragraph{Designers Can Execute a Simulated Run}
\textbf{Value:} 2\quad \textbf{Priority:} Should\\
\textbf{Acceptance criteria}
\begin{itemize}
  \item A designer is able to execute a simulated run.
  \item The run uses a player simulation instead of a user player, and enemies behave as normal.
  \item The player agent is able to perform any action a user player could.
  \item The run continues until the agent beats the game or dies.
\end{itemize}

\paragraph{Developers Can Assign Roles}
\textbf{Value:} 5\quad \textbf{Priority:} Could\\
\textbf{Acceptance criteria}
\begin{itemize}
  \item An authenticated developer can enter the username of another user and assign a role to them.
  \item The first user must be a developer.
  \item There must always be one developer account.
\end{itemize}

\subsubsection{Game Runs}

\paragraph{User Can Choose Difficulty}
\textbf{Value:} 3\quad \textbf{Priority:} Must\\
\textbf{Acceptance criteria}
\begin{itemize}
  \item At the start of a game run, the player can select the difficulty for the run.
  \item The difficulty options are easy, medium, and hard.
  \item The game must be harder on each subsequent difficulty option (easy is the easiest, hard is the hardest).
  \item The game must vary how hard a difficulty is by varying the values of design parameters.
\end{itemize}

\paragraph{Player Can Use Mundane Attack Abilities}
\textbf{Value:} 8\quad \textbf{Priority:} Must\\
\textbf{Acceptance criteria}
\begin{itemize}
  \item The player must always be able to use the default attack ability on their turn while in an encounter.
  \item The player must be able to use any mundane attack abilities they have purchased in the shop on their turn in an encounter.
  \item These abilities must deal the damage type specified in the game design documentation.
  \item These abilities must deal an amount of damage based on the game design documentation to the enemy.
  \item The enemy's health must be reduced by this damage amount.
  \item Any purchasable mundane attack abilities must be purchasable in the shop between encounters.
  \item They must not be purchasable if the player does not have enough coins.
  \item They must reduce the amount of coins the player has by their cost when purchased.
  \item They must only be purchased once per run.
\end{itemize}

\paragraph{Player Can Choose Which Abilities to Purchase}
\textbf{Value:} 1\quad \textbf{Priority:} Should\\
\textbf{Acceptance criteria}
\begin{itemize}
  \item Multiple abilities must appear in the shop between encounters.
  \item Players must be able to choose which ability from the shop they wish to buy.
  \item Players must be able to purchase multiple abilities from the same shop.
\end{itemize}

\paragraph{Player Can Use Purchased Magic Attack Abilities}
\textbf{Value:} 8\quad \textbf{Priority:} Could\\
\textbf{Acceptance criteria}
\begin{itemize}
  \item Players must have a finite magic reserve, limited by their maximum magic.
  \item While in an encounter, if the player has sufficient magic, they must be able to use a magic ability on their turn if they have purchased it.
  \item If the player does not have enough magic to fuel a magic attack ability, they must not be able to use it.
  \item When a magic attack ability is used, it must reduce their magic by its cost as defined in the game design documentation.
  \item These abilities must deal the damage type specified in the game design documentation.
  \item These abilities must deal an amount of damage based on the game design documentation to the enemy.
  \item The enemy's health must be reduced by this damage amount.
  \item Any purchasable magic attack abilities must be purchasable in the shop between encounters.
  \item They must not be purchasable if the player does not have enough coins.
  \item They must reduce the amount of coins the player has by their cost when purchased.
  \item They must only be purchased once per run.
\end{itemize}

\paragraph{Enemies Can Use Attack Abilities}
\textbf{Value:} 3\quad \textbf{Priority:} Must\\
\textbf{Acceptance criteria}
\begin{itemize}
  \item Enemies must be able to use both magic and mundane attack abilities on their turn against the player.
  \item Enemies must have the attack abilities based on their enemy type defined in the game design documentation.
  \item These abilities must deal damage of the type defined in the game design documentation.
  \item These abilities must deal an amount of damage based on the game design documentation, to the player.
  \item The player's health must be reduced by this amount of damage.
  \item An enemy whose health is zero or less is dead and must not be able to use attack abilities.
\end{itemize}

\paragraph{Player Benefits From Purchased Damage Resistance Passive Abilities}
\textbf{Value:} 2\quad \textbf{Priority:} Could\\
\textbf{Acceptance criteria}
\begin{itemize}
  \item The player must be able to purchase passive abilities in the shop that provide resistance to a specific damage type, as specified in the game design documentation.
  \item These abilities must affect the player in an encounter by halving the damage they take when the damage being dealt to them is the type that the upgrade provides resistance to.
  \item These must not be purchasable if the player does not have enough coins.
  \item These must reduce the amount of coins the player has by their cost when purchased.
  \item They must only be purchased once per run.
\end{itemize}

\paragraph{Player Benefits From Purchased Damage Boost Passive Abilities}
\textbf{Value:} 2\quad \textbf{Priority:} Could\\
\textbf{Acceptance criteria}
\begin{itemize}
  \item The player must be able to purchase passive abilities in the shop that provide boosts to a specific damage type, as specified in the game design documentation.
  \item These abilities must affect the player in an encounter by doubling the damage they deal when the damage being dealt by them is the type that the upgrade boosts.
  \item These must not be purchasable if the player does not have enough coins.
  \item These must reduce the amount of coins the player has by their cost when purchased.
  \item They must only be purchased once per run.
\end{itemize}

\paragraph{Enemies Have Damage Resistances and Vulnerabilities}
\textbf{Value:} 2\quad \textbf{Priority:} Could\\
\textbf{Acceptance criteria}
\begin{itemize}
  \item Some enemies must resist a damage type as specified in the game design documentation.
  \item Some enemies must be weak to a damage type as specified in the game design documentation.
  \item When an enemy is dealt damage of a type they resist, that damage must be halved.
  \item When an enemy is dealt damage of a type they are vulnerable to, that damage must be doubled.
\end{itemize}

\paragraph{Players Can Encounter Multiple Enemies in Combat}
\textbf{Value:} 2\quad \textbf{Priority:} Could\\
\textbf{Acceptance criteria}
\begin{itemize}
  \item Combat encounters must be able to have multiple enemies within them.
  \item Each enemy must target the player with their attack abilities.
  \item The player must be able to choose which enemy to target with their attack abilities.
\end{itemize}

\paragraph{Players Can Complete a Combat Round}
\textbf{Value:} 1\quad \textbf{Priority:} Must\\
\textbf{Acceptance criteria}
\begin{itemize}
  \item At the start of a combat round, the player must gain magic based on the game design specification.
  \item They must then take their turn.
  \item All enemies must then take their turns.
  \item The round is then complete and starts from the beginning, granting the player magic.
\end{itemize}

\paragraph{Players Can Complete an Encounter}
\textbf{Value:} 1\quad \textbf{Priority:} Must\\
\textbf{Acceptance criteria}
\begin{itemize}
  \item A combat encounter must be completed if all enemies within the encounter have zero or less health.
  \item When complete, the player must gain coins based on the game design documentation.
\end{itemize}

\paragraph{Players Can Retry an Encounter}
\textbf{Value:} 1\quad \textbf{Priority:} Must\\
\textbf{Acceptance criteria}
\begin{itemize}
  \item At the start of a game run, the player must have an amount of lives based on the design parameter.
  \item If the player's health reaches zero or less in an encounter, they must die.
  \item If a player dies while having one or more lives, they must restart the same encounter. Their health and magic must be reset to their starting values, as must the health of all enemies within the encounter.
  \item If a player dies, their lives must be reduced by one.
\end{itemize}

\paragraph{Players Can Progress Through Stages}
\textbf{Value:} 1\quad \textbf{Priority:} Must\\
\textbf{Acceptance criteria}
\begin{itemize}
  \item After completing an encounter, the player must go to the shop.
  \item After leaving the shop, the stage the player is on must be incremented.
  \item The player must then enter another encounter.
\end{itemize}

\paragraph{Players Can Fail Runs}
\textbf{Value:} 1\quad \textbf{Priority:} Must\\
\textbf{Acceptance criteria}
\begin{itemize}
  \item If a player dies during an encounter while having zero or fewer lives, they must fail the run.
  \item If a player fails the run, the run must end, displaying a run-over screen.
  \item The screen must show what stage the player reached, how long the run was, how many coins they had, and how many times they died.
  \item The player must lose all their purchased abilities and have their stage reset along with all encounters.
  \item The player must then be able to start a new run after leaving this screen.
\end{itemize}

\paragraph{Players Can Complete Runs}
\textbf{Value:} 1\quad \textbf{Priority:} Must\\
\textbf{Acceptance criteria}
\begin{itemize}
  \item If a player completes the encounter on stage two, they must complete the run.
  \item This shows a run completed screen displaying what stage the player reached, how long the run was, how many coins they finished with, and how many times they died.
  \item They must then be able to leave this screen and start a new run.
\end{itemize}

\subsubsection{Telemetry Events}

\paragraph{Application Handles Start Session}
\textbf{Value:} 1\quad \textbf{Priority:} Must\\
\textbf{Acceptance criteria}
\begin{itemize}
  \item The game features the start session event.
  \item The event has the fields specified in the telemetry schema.
  \item The event is created when a new run starts.
\end{itemize}

\paragraph{Application Handles Normal Encounter Start}
\textbf{Value:} 1\quad \textbf{Priority:} Must\\
\textbf{Acceptance criteria}
\begin{itemize}
  \item The game features the normal encounter start event.
  \item The event has the fields specified in the telemetry schema.
  \item The event is created each time a normal encounter starts, including restarts.
\end{itemize}

\paragraph{Application Handles Normal Encounter Fail}
\textbf{Value:} 1\quad \textbf{Priority:} Must\\
\textbf{Acceptance criteria}
\begin{itemize}
  \item The game features the normal encounter fail event.
  \item The event has the fields specified in the telemetry schema.
  \item The event is created each time a normal encounter is failed by the user's character dying.
\end{itemize}

\paragraph{Application Handles Normal Encounter Complete}
\textbf{Value:} 1\quad \textbf{Priority:} Must\\
\textbf{Acceptance criteria}
\begin{itemize}
  \item The game features the normal encounter complete event.
  \item The event has the fields specified in the telemetry schema.
  \item The event is created each time the player completes a normal encounter by killing all the enemies.
\end{itemize}

\paragraph{Application Handles Boss Encounter Start}
\textbf{Value:} 0\quad \textbf{Priority:} Wont\\
\textbf{Acceptance criteria} N/A

\paragraph{Application Handles Boss Encounter Fail}
\textbf{Value:} 0\quad \textbf{Priority:} Wont\\
\textbf{Acceptance criteria} N/A

\paragraph{Application Handles Boss Encounter Complete}
\textbf{Value:} 0\quad \textbf{Priority:} Wont\\
\textbf{Acceptance criteria} N/A

\paragraph{Application Handles Gain Coin}
\textbf{Value:} 1\quad \textbf{Priority:} Could \\
\textbf{Acceptance criteria}
\begin{itemize}
  \item The game features the gain coin event.
  \item The event has the fields specified in the telemetry schema.
  \item The event is created when the player gains coins.
\end{itemize}

\paragraph{Application Handles Buy Upgrade}
\textbf{Value:} 1\quad \textbf{Priority:} Could \\
\textbf{Acceptance criteria}
\begin{itemize}
  \item The game features the buy upgrade event.
  \item The event has the fields specified in the telemetry schema.
  \item The event is created when the player buys an ability.
\end{itemize}

\paragraph{Application Handles End Session}
\textbf{Value:} 1\quad \textbf{Priority:} Must\\
\textbf{Acceptance criteria}
\begin{itemize}
  \item The game features the end session event.
  \item The event has the fields specified in the telemetry schema.
  \item The event is created when a game run ends, either by the user completing the run by completing the second encounter, or by their character dying.
\end{itemize}

\paragraph{Application Handles Settings Change}
\textbf{Value:} 1\quad \textbf{Priority:} Must\\
\textbf{Acceptance criteria}
\begin{itemize}
  \item The game features the settings change event.
  \item The event has the fields specified in the telemetry schema.
  \item The event is created when the user changes a setting's value and saves this change.
\end{itemize}

\paragraph{Application Handles Kill Enemy}
\textbf{Value:} 1\quad \textbf{Priority:} Could \\
\textbf{Acceptance criteria}
\begin{itemize}
  \item The game features the kill enemy event.
  \item The event has the fields specified in the telemetry schema.
  \item The event is created when the player kills an enemy.
\end{itemize}
